\documentclass[11pt]{article}
\usepackage[a4paper,margin=2.4cm]{geometry}
\usepackage{amsmath,amssymb,bm}
\usepackage{slashed}
\usepackage{booktabs}
\usepackage{array}
\usepackage{hyperref}
\usepackage{xcolor}
\usepackage{listings}

\hypersetup{colorlinks=true,linkcolor=blue,citecolor=blue,urlcolor=blue}
\lstset{
  basicstyle=\ttfamily\small,
  breaklines=true,
  frame=single,
  columns=fullflexible
}

\newcommand{\as}{\alpha_s}
\newcommand{\MSbar}{\overline{\mathrm{MS}}}
\newcommand{\slashedp}{\slashed{p}}

\title{Independent \(O(\alpha_s)\) Calculation for the \(B_c\) Axial Current\\
Lecture Notes, Validation Tests, and Current Audit Status}
\author{}
\date{\today}

\begin{document}
\maketitle

\begin{abstract}
These notes describe the momentum-space calculation of the perturbative
\(O(\alpha_s)\) correction to the axial-current correlator used in the
\(B_c\) mixing analysis.  The purpose is twofold: to give a derivation that a
graduate student can reproduce and to expose every convention needed for an
independent theoretical check.  The leading-order, soft, ultraviolet-pole,
infrared-pole, equal-mass form-factor, scale-cancellation, and small-mass
checks pass.  Two normalization errors in an earlier provisional result are
identified explicitly: the Package-X epsilon-bar symbol was inverted, and
the \(O(\epsilon)\) difference between the \(D\)-dimensional and
four-dimensional projected Born traces was omitted.  After correcting both,
the independent axial coefficient reproduces the known equal-mass result and
has a smooth massless limit.  Numerical results are quoted separately in an
on-shell diagnostic convention and in a common-scale \(\MSbar\) convention.
\end{abstract}

\tableofcontents

\section{Physics goal and conventions}

The axial current is
\begin{equation}
 J_\mu^A(x)=\bar b(x)\gamma_\mu\gamma_5 c(x),
\end{equation}
and the two-point function is
\begin{equation}
 \Pi_{\mu\nu}^{AA}(p)
 =i\int d^4x\,e^{ip\cdot x}
 \langle0|T\{J_\mu^A(x)J_\nu^{A\dagger}(0)\}|0\rangle .
\end{equation}
We decompose
\begin{equation}
 \Pi_{\mu\nu}^{AA}(p)
 =\left(g_{\mu\nu}-\frac{p_\mu p_\nu}{p^2}\right)\Pi_1^{AA}(p^2)
 +\frac{p_\mu p_\nu}{p^2}\Pi_0^{AA}(p^2)
\end{equation}
and project with
\begin{equation}
 P_1^{\mu\nu}
 =\frac{1}{3}\left(g^{\mu\nu}-\frac{p^\mu p^\nu}{p^2}\right)
\end{equation}
in four dimensions.  In the virtual loop, the projector must instead be
\begin{equation}
 P_{1,D}^{\mu\nu}
 =\frac{1}{D-1}\left(g_D^{\mu\nu}
 -\frac{p^\mu p^\nu}{p^2}\right),\qquad D=4-2\epsilon .
\end{equation}

The perturbative spectral density is defined by
\begin{equation}
 \rho_{\rm pert}^{AA}(s)
 =\rho_0^{AA}(s)+\frac{\as}{\pi}\rho_1^{AA}(s)+O(\as^2).
\end{equation}
Thus \(\rho_1\) never contains the explicit factor \(\as/\pi\).

\section{Leading-order cut}

For the cut state \(c(p_c)\bar b(p_b)\), \(p=p_c+p_b\), the spin sum is
\begin{equation}
 T_{\mu\nu}^{(0)}
 =\mathrm{Tr}\left[
 (\slashed p_c+m_c)\gamma_\mu\gamma_5
 (\slashed p_b-m_b)\gamma_\nu\gamma_5
 \right].
\end{equation}
The on-shell relations are
\begin{equation}
 p_c^2=m_c^2,\qquad p_b^2=m_b^2,\qquad
 2p_c\cdot p_b=s-m_c^2-m_b^2.
\end{equation}
The two-body phase-space factor is proportional to
\begin{equation}
 \frac{\sqrt{\lambda(s,m_b^2,m_c^2)}}{s},\qquad
 \lambda(s,m_b^2,m_c^2)
 =[s-(m_b+m_c)^2][s-(m_b-m_c)^2].
\end{equation}
In the normalization of the momentum-space mixing code,
\begin{equation}
 \rho_0^{AA}(s)
 =\frac{N_c}{16\pi^2}\frac{\sqrt{\lambda}}{s}
 \,\mathcal N_{AA}(s).
\end{equation}
The independently evaluated trace agrees algebraically with
\texttt{PerturbativeSpectralDensity["AA",s]}:
\begin{lstlisting}[language=Mathematica]
IndependentAALODerivationCheck[s]["Difference"]
(* 0 *)
\end{lstlisting}

\section{NLO decomposition}

The inclusive correction is divided into
\begin{equation}
 \rho_1^{AA}
 =\rho_{1,V}^{AA}
 +\rho_{1,R-D}^{AA}
 +\rho_{1,I}^{AA},
\end{equation}
where \(V\) is the renormalized virtual correction, \(R-D\) is the exact
three-body real contribution minus local dipoles, and \(I\) is the integrated
dipole insertion.  The individual terms depend on subtraction conventions;
only their sum is physical.

\section{Virtual correction}

\subsection{Loop amplitude}

With loop momentum \(\ell\), the vertex denominators are represented by
\begin{equation}
 \frac{1}{
 \ell^2\big[(\ell+p_c)^2-m_c^2\big]
 \big[(\ell-p_b)^2-m_b^2\big]} .
\end{equation}
The numerator is traced in \(D\) dimensions.  In FeynCalc this requires
\texttt{GAD}, \texttt{GSD}, \texttt{MTD}, and \texttt{FVD}; performing a
four-dimensional trace and applying \texttt{ChangeDimension} afterward is
not equivalent because \(O(\epsilon)\) numerator terms can multiply
\(1/\epsilon\) poles.

There are two nonsinglet \(\gamma_5\) matrices.  They are anticommuted
together before the trace is continued to \(D\) dimensions.  Thus no
four-dimensional \(\gamma_5\) object is mixed with a \(D\)-dimensional gamma
chain.  The corresponding four-dimensional identity is checked by
\begin{lstlisting}[language=Mathematica]
IndependentAALOGamma5EliminationCheck[]
(* 0 *)
\end{lstlisting}

The tensor integrals are reduced with
\begin{equation}
 \texttt{TID}\longrightarrow\texttt{ToPaVe}
 \longrightarrow\texttt{PaXEvaluateUVIRSplit}.
\end{equation}
The physical loop measure \(d^D\ell/(2\pi)^D\) is supplied through
\texttt{PaXImplicitPrefactor}.

\subsection{Epsilon-bar convention and projected Born normalization}

FeynHelpers defines
\begin{equation}
 \frac{1}{\texttt{PaXEpsilonBar}}
 =\frac{1}{\epsilon}-\gamma_E+\ln(4\pi).
\end{equation}
Therefore the Laurent-series substitution is
\(\texttt{PaXEpsilonBar}\to z\), not \(1/z\).  Reversing this substitution
was the dominant error in the superseded \(+3\%\) result.

The second necessary correction follows from the projected Born trace.  The
ratio between the \(D\)-dimensional gamma-five-free trace and the
four-dimensional trace is
\begin{equation}
 \frac{B_D}{B_4}
 =1+\epsilon\,
 \frac{2[(m_b-m_c)^2-s]}
 {3(m_b-m_c)^2+6s}+O(\epsilon^2).
\end{equation}
For equal masses this becomes \(1-\epsilon/3\).  Since the virtual graph
contains \(1/\epsilon\) poles, this apparently evanescent factor changes the
finite result.  The loop is consequently multiplied by \(B_4/B_D\) before
the finite part is extracted.

\subsection{On-shell field counterterm}

For the coefficient of \(\as/(4\pi)\),
\begin{equation}
 \delta Z_{2,Q}^{\rm OS}
 =-C_F\left[
 \frac{1}{\epsilon_{\rm UV}}
 +\frac{2}{\epsilon_{\rm IR}}
 +4+3\ln\frac{\mu^2}{m_Q^2}
 \right].
\end{equation}
The interference of the two external-field factors with the Born amplitude
gives
\begin{equation}
 \rho_{1,\rm field}^{AA}
 =\frac14\left(\delta Z_{2,b}^{\rm OS}
 +\delta Z_{2,c}^{\rm OS}\right)\rho_0^{AA}.
\end{equation}
The factor \(1/4\) converts the \(\as/(4\pi)\) convention into the
\(\as/\pi\) convention.

\section{Real-gluon emission}

For
\begin{equation}
 J_A\longrightarrow c(p_c)+\bar b(p_b)+g(k),
\end{equation}
define
\begin{equation}
 u=2p_c\cdot k,\qquad v=2p_b\cdot k,\qquad
 t=(p_c+p_b)^2=s-u-v.
\end{equation}
The two emission chains are
\begin{align}
 {\cal M}_c^\alpha&=
 \gamma^\alpha
 \frac{\slashed p_c+\slashed k+m_c}{u}
 \gamma_\mu\gamma_5,\\
 {\cal M}_{\bar b}^\alpha&=
 \gamma_\mu\gamma_5
 \frac{-\slashed p_b-\slashed k+m_b}{v}
 \gamma^\alpha.
\end{align}
Their sum obeys the soft theorem.  The three-body measure after integrating
the overall orientation is
\begin{equation}
 d\Phi_3=\frac{dt\,du}{128\pi^3s}.
\end{equation}
The spectral density includes the optical-theorem factor \(1/(2\pi)\).

At fixed \(t\),
\begin{equation}
 u_\pm(t)=\frac{s-t}{2t}
 \left[t+m_c^2-m_b^2
 \pm\sqrt{\lambda(t,m_b^2,m_c^2)}\right].
\end{equation}

\section{Local massive dipoles}

For a massive emitter \(Q\), emitted gluon \(g\), and massive spectator \(k\),
the Catani--Dittmaier--Seymour--Trocsanyi dipole has the form
\begin{equation}
 D_{gQ,k}
 =\frac{1}{2p_g\cdot p_Q}
 \langle V_{gQ,k}\rangle |{\cal M}_0(\widetilde p_Q,\widetilde p_k)|^2 .
\end{equation}
The local subtraction was checked analytically:
\begin{lstlisting}[language=Mathematica]
IndependentAADipoleSoftCancellationCheck[s,yAA]["CancelsQ"]
(* True *)
\end{lstlisting}
This is necessary but not sufficient: the quasi-collinear and integrated
finite limits must also be checked.

\section{Complete integrated insertion}

A key lesson from the audit is that the individual integrated dipole
\(I_{gQ,k}\) is not by itself the complete finite insertion.  For final-state
partons, the insertion operator is
\begin{align}
 {\bf I}_m(\epsilon)
 =-&\frac{\as}{2\pi}\frac{(4\pi)^\epsilon}{\Gamma(1-\epsilon)}
 \sum_j\frac{1}{{\bf T}_j^2}\sum_{k\ne j}{\bf T}_j\cdot{\bf T}_k
 \Bigg[
 {\bf T}_j^2\left(\frac{\mu^2}{s_{jk}}\right)^\epsilon
 \left({\cal V}_j-\frac{\pi^2}{3}\right)\nonumber\\
 &+\Gamma_j+\gamma_j\ln\frac{\mu^2}{s_{jk}}
 +\gamma_j+K_j
 \Bigg].
\end{align}
For a quark,
\begin{equation}
 \gamma_q=\frac32C_F,\qquad
 K_q=\left(\frac72-\frac{\pi^2}{6}\right)C_F,
\end{equation}
and for a massive quark,
\begin{equation}
 \Gamma_q(\mu,m_q;\epsilon)
 =C_F\left[
 \frac1\epsilon+\frac12\ln\frac{m_q^2}{\mu^2}-2
 \right].
\end{equation}
The functions \({\cal V}^{(S)}\) and \({\cal V}_q^{(NS)}\) are implemented
from Eqs.\ (6.20) and (6.21) of
Catani et al., Nucl.\ Phys.\ B627 (2002) 189.

\section{Pole and scale checks}

The following tests currently pass:
\begin{enumerate}
 \item The LO cut reproduces the stored spectral density.
 \item The exact real matrix element reproduces the universal massive soft
 eikonal factor.
 \item The virtual UV pole cancels the on-shell field counterterm.
 \item The real and virtual IR-pole coefficients cancel.
 \item The derivative of the finite virtual term with respect to
 \(\ln\mu^2\) equals its IR-pole coefficient.
\end{enumerate}

The scale-slope test is particularly useful because it checks the dimensional
prefactor without depending on a published full spectral density.

\section{How to validate a large correction}

A large correction is not disproved by its size alone.  Near a heavy-heavy
threshold, gluon exchange is enhanced by the small relative velocity.
Nevertheless, a large result requires more tests than a small one.  It is
useful to organize them in three levels.

\subsection{Level I: algebraic and numerical identities}

These tests use no external spectral-density formula:
\begin{enumerate}
 \item exchange the two masses, \(m_b\leftrightarrow m_c\);
 \item rescale every dimensional variable;
 \item change the numerical quadrature order;
 \item vary the subtraction cutoff or integration variables;
 \item verify cancellation of the renormalization scale in the complete
 on-shell coefficient.
\end{enumerate}

The \(AA\) density has mass dimension two.  Therefore
\begin{equation}
 \rho_1(\kappa^2s,\kappa m_b,\kappa m_c,\kappa\mu)
 =\kappa^2\rho_1(s,m_b,m_c,\mu).
\end{equation}
For \(\kappa=2\), the numerical ratio between the left-hand side and the
right-hand side is \(0.9999988\).  Exchanging \(m_b\) and \(m_c\) changes
\(\rho_1(s=40\,{\rm GeV}^2)\) by only
\(1.8\times10^{-14}\).

\subsection{Level II: limits derived directly from QCD}

\paragraph{Threshold or NRQCD limit.}
For two heavy quarks the appropriate effective description near threshold
is NRQCD rather than ordinary heavy-light HQET.  Define
\begin{equation}
 v_{\rm rel}
 =\frac{\sqrt{\lambda(s,m_b^2,m_c^2)}}
 {s-m_b^2-m_c^2}.
\end{equation}
Single Coulomb-gluon exchange implies
\begin{equation}
 \frac{\rho_1}{\rho_0}
 =\frac{C_F\pi^2}{v_{\rm rel}}+O(\ln v_{\rm rel},1).
\end{equation}
This is the first term in the expansion of the Sommerfeld/Coulomb factor.
For the physical mass inputs the numerical check is
\begin{center}
\begin{tabular}{cccc}
\toprule
\(s-s_{\rm th}\,[{\rm GeV}^2]\) & \(v_{\rm rel}\) &
\(\rho_1/\rho_0\) & \(v_{\rm rel}\rho_1/\rho_0\)\\
\midrule
0.10 & 0.13629 & 95.1736 & 12.9711\\
0.25 & 0.21326 & 60.1795 & 12.8340\\
0.50 & 0.29653 & 42.6488 & 12.6465\\
1.00 & 0.40589 & 30.3493 & 12.3186\\
\bottomrule
\end{tabular}
\end{center}
The limiting expectation is \(C_F\pi^2=13.1595\).  The approach toward this
number is a strong explanation for the large Borel correction: the Borel
weight emphasizes the near-threshold region.  It also warns that fixed-order
perturbation theory may eventually require Coulomb resummation.

\paragraph{High-energy limit.}
For \(s\gg m_b^2,m_c^2\), both masses become negligible and the nonsinglet
axial result approaches the massless value
\begin{equation}
 \frac{\rho_1}{\rho_0}\longrightarrow1.
\end{equation}
For the physical masses, the calculated ratios are
\begin{equation}
\begin{array}{c|ccc}
s\,[{\rm GeV}^2]&100&400&1600\\ \hline
\rho_1/\rho_0&4.0352&1.9458&1.3035
\end{array}.
\end{equation}
The convergence is slow because the bottom mass is not negligible at the
first two points, but the sequence approaches one.

\paragraph{Heavy-light limit.}
Sending one mass to zero while keeping the other fixed provides a different
check.  Ordinary HQET becomes relevant near the heavy-light threshold.  At
\(s=40\,{\rm GeV}^2\), \(m_b=4.18\,{\rm GeV}\), and
\(m_c=1.0,0.5,0.2,0.1,0.05\,{\rm GeV}\), the ratios
\(\rho_1/\rho_0\) are
\begin{equation}
 8.4611,\quad6.4876,\quad5.6258,\quad5.3519,\quad5.2036.
\end{equation}
The sequence is finite and smooth.  A stronger check is to match the exact
normalization to the known heavy-light correlator of
Chetyrkin and Steinhauser~\cite{ChetyrkinSteinhauser2001}.  Because their
published basis is not identical to the present projected axial convention,
that comparison must be performed at the level of a carefully matched
spectral density rather than by comparing one number.

\subsection{Level III: independent published calculations}

The following comparisons probe different parts of the result:
\begin{enumerate}
 \item Catani et al.~\cite{Catani2002}: massive subtraction and the exact
 equal-mass virtual form factors.
 \item Chetyrkin and Steinhauser~\cite{ChetyrkinSteinhauser2001}: the
 heavy-light, one-massless-quark limit and its HQET threshold expansion.
 \item Lee, Sang and Kim~\cite{LeeSangKim2010}: unequal-mass
 \(b\bar c\) current matching in NRQCD, including the one-loop and
 relativistic expansion.  Their principal application is to S-wave
 \(B_c\) currents, whereas the transverse axial channel studied here has a
 P-wave threshold structure.  The comparison is therefore limited to
 carefully matched Lorentz components and short-distance conventions; it is
 not a direct spectral-density check.
 \item Maier and Marquard~\cite{MaierMarquard2011}: high-energy expansions
 of non-diagonal vector/scalar correlators and diagonal axial correlators.
 These results test the expected expansion structure, but they are not a
 published unequal-mass axial formula for the present current.
 \item Wang~\cite{Wang2013}: a direct unequal-mass NLO spectral-density
 comparison, treated only as an external check because the local
 transcription still requires an independent convention audit.
\end{enumerate}

\subsection{What the old \(K\)-factor model checks}

The earlier diagnostic replaced
\begin{equation}
 \Pi_{\rm pert}^{ij}\to
 \left(1+\frac{\as}{\pi}K_{ij}\right)\Pi_{\rm pert}^{ij}.
\end{equation}
It remains useful for propagating hypothetical channel-dependent corrections
through
\begin{equation}
 \tan(2\theta)=
 \frac{-2\Pi^{AB}}{\Pi^{AA}-\Pi^{BB}}.
\end{equation}
If all \(K_{ij}\) are equal, a common correction largely cancels in the
angle.  If they differ, the angle can move by several degrees.  The scan
therefore measures the \emph{sensitivity of the angle}; it cannot validate
\(\rho_1^{AA}(s)\), reproduce threshold logarithms, or assign a statistical
NLO uncertainty.

\section{Mandatory limiting tests}

\subsection{Massless limit}

For equal masses followed by \(m^2/s\to0\), the transverse nonsinglet axial
current becomes the massless vector/axial current.  Therefore
\begin{equation}
 \frac{\rho_1^{AA}(s)}{\rho_0^{AA}(s)}\longrightarrow1.
\end{equation}
At \(s=40\,{\rm GeV}^2\), the corrected numerical sequence is
\begin{center}
\begin{tabular}{ccc}
\toprule
\(m_b=m_c\) [GeV] & \(\rho_1^{AA}/\rho_0^{AA}\)\\
\midrule
0.20 & 1.0831974\\
0.10 & 1.0249053\\
0.05 & 1.0072604\\
\bottomrule
\end{tabular}
\end{center}
which approaches one as required.

\subsection{Equal-mass form factor}

The virtual contribution is checked against the equal-mass form factors in
Appendix D of Catani et al.  Both \(f_1\) and the finite-mass pure-axial
\(f_2\) term are included.  At \(s=40\,{\rm GeV}^2\) and
\(\mu=\sqrt{s}\), the comparison is
\begin{center}
\begin{tabular}{cccc}
\toprule
\(m\) [GeV] & FeynCalc/Package-X & Catani axial & difference\\
\midrule
0.20 & 42.54933479185 & 42.54933479167 & \(1.75\times10^{-10}\)\\
0.50 & 26.78426097716 & 26.78426097716 & \(-4.71\times10^{-12}\)\\
1.00 & 18.02429851421 & 18.02429851421 & \(2.49\times10^{-14}\)\\
\bottomrule
\end{tabular}
\end{center}
This validates the finite virtual normalization independently of the
unequal-mass literature.

\subsection{Unequal-mass external comparison}

The expressions of Z.-G. Wang, arXiv:1303.4146, may be used only after:
\begin{enumerate}
 \item reproducing their LO normalization;
 \item checking their equal-mass limit analytically rather than substituting
 directly into formulas containing \((m_1-m_2)^{-1}\);
 \item checking the massless limit;
 \item comparing several values of \(s\), not one integrated number.
\end{enumerate}
The current local transcription of that formula does not pass all these
requirements and is not used to construct the result.

\section{Borel transformation}

Once a validated \(\rho_1^{AA}\) is available, its continuum-subtracted
Borel moment is
\begin{equation}
\Pi_{1,\rm NLO}^{AA}(M^2,s_0)
 =\frac{\as(\mu)}{\pi}
 \int_{(m_b+m_c)^2}^{s_0}ds\,
 e^{-s/M^2}\rho_1^{AA}(s,\mu).
\end{equation}
For the fixed numerical inputs \(m_b=4.18\,{\rm GeV}\),
\(m_c=1.27\,{\rm GeV}\), \(\as=0.26\),
\(M^2=10\,{\rm GeV}^2\), and \(s_0=55\,{\rm GeV}^2\), interpreted only as an
on-shell-scheme diagnostic, the pointwise decomposition at
\(s=40\,{\rm GeV}^2\) is
\begin{equation}
 \rho_{1,V}=5.6979732,\qquad
 \rho_{1,I}=-2.2122115,\qquad
 \rho_{1,R-D}=0.4477501,
\end{equation}
so that \(\rho_1^{AA}=3.9335117\).  The LO density is
\(\rho_0^{AA}=0.3898262\), and the local relative correction
\((\as/\pi)\rho_1/\rho_0\) is \(0.8351\).

The continuum-subtracted 20-node Borel result is
\begin{align}
 \Pi_{\rm LO}^{AA}&=0.1441716826,\\
 \Pi_{1}^{AA}&=1.3770512496,\\
 \frac{\as}{\pi}\Pi_{1}^{AA}&=0.1139655469,\\
 \Pi_{\rm LO+NLO}^{AA}&=0.2581372295.
\end{align}
Thus the relative correction is \(0.790485\), or \(79.05\%\).  The 12-node
and 20-node bare coefficients are \(1.3770552219\) and \(1.3770512496\),
respectively.  This agreement tests the outer Borel quadrature.

This large correction is not combined with the paper's
\(\MSbar\) mass inputs, because \(4.18\) and \(1.27\) GeV are normally
interpreted as \(m_b(m_b)\) and \(m_c(m_c)\), not pole masses.

\section{\(\MSbar\) conversion}

Conversion to running masses uses
\begin{equation}
 M_Q=\overline m_Q(\mu)\left[
 1+\frac{\as(\mu)}{\pi}
 \left(\frac43+\ln\frac{\mu^2}{\overline m_Q^2(\mu)}\right)
 \right],
\end{equation}
which implies
\begin{equation}
 \rho_{1,\MSbar}^{AA}
 =\rho_{1,\rm OS}^{AA}
 +\sum_{Q=b,c}\overline m_Q(\mu)
 \left(\frac43+\ln\frac{\mu^2}{\overline m_Q^2(\mu)}\right)
 \frac{\partial\rho_0^{AA}}{\partial m_Q}.
\end{equation}

Using \(m_b(m_b)=4.18\,{\rm GeV}\), \(m_c(m_c)=1.27\,{\rm GeV}\),
\(\as(M_Z)=0.1180\), one-loop running, and the common scale
\(\mu=2.0\,{\rm GeV}\), the parameters are
\begin{equation}
 \overline m_b(\mu)=4.6683009\,{\rm GeV},\qquad
 \overline m_c(\mu)=1.1659209\,{\rm GeV},\qquad
 \as(\mu)=0.2622100.
\end{equation}
The corrected 20-node result is
\begin{align}
 \Pi_{\rm LO,\MSbar}^{AA}&=0.0723406411,\\
 \Pi_{1,\rm OS}^{AA}\big|_{\overline m(\mu)}&=0.7641393212,\\
 \Pi_{1,\rm conversion}^{AA}&=-0.1431034913,\\
 \Pi_{1,\MSbar}^{AA}&=0.6210358299,\\
 \frac{\as(\mu)}{\pi}\Pi_{1,\MSbar}^{AA}&=0.0518341572,\\
 \Pi_{\rm LO+NLO,\MSbar}^{AA}&=0.1241747983.
\end{align}
The relative NLO correction is \(71.65\%\).  Varying the common scale gives
\begin{center}
\begin{tabular}{cccc}
\toprule
\(\mu\) [GeV] & \(\Pi_{\rm LO}^{AA}\) &
\((\as/\pi)\Pi_1^{AA}\) & \(\Pi_{\rm LO+NLO}^{AA}\)\\
\midrule
1.8 & 0.05866679 & 0.05981541 & 0.11848220\\
2.0 & 0.07234064 & 0.05183386 & 0.12417450\\
2.2 & 0.08587655 & 0.04205294 & 0.12792949\\
\bottomrule
\end{tabular}
\end{center}
The central value should therefore be quoted with a renormalization-scale
band.  The fixed-input on-shell and common-scale \(\MSbar\) numbers answer
different scheme questions and must not be added to each other.

\section{Reproducible Mathematica audit}

\begin{lstlisting}[language=Mathematica]
Get["BcMixingAlphaS.wl"];

IndependentAALODerivationCheck[s]["Difference"]
IndependentAARealEmissionSoftTheoremCheck[s,yAA]["Difference"]
IndependentAAUVPoleCancellationCheck[
  s,l,FeynCalc`PaXC0Expand->True
]["CancelsQ"]
IndependentAAIRPoleCancellationCheck[40.0][
  "CancelsNumericallyQ"
]

IndependentAAEqualMassVirtualBenchmarkReport[
  40,1/5,Sqrt[40]
]["PassesFiniteBenchmarkQ"]

IndependentAAThresholdCoulombCheck[
  {0.1,0.25,0.5,1.0},
  $BcMixingDefaultParameters,
  AccuracyGoal->5,PrecisionGoal->5,MaxRecursion->10
]

IndependentAAHighEnergyLimitCheck[
  {100.,400.,1600.},
  $BcMixingDefaultParameters,
  AccuracyGoal->5,PrecisionGoal->5,MaxRecursion->10
]

IndependentAAMassExchangeSymmetryCheck[
  40.,$BcMixingDefaultParameters,
  AccuracyGoal->5,PrecisionGoal->5,MaxRecursion->10
]

IndependentAADimensionalScalingCheck[
  40.,2.,$BcMixingDefaultParameters,
  AccuracyGoal->5,PrecisionGoal->5,MaxRecursion->10
]

aaOS = IndependentAAFinalNLOBorelSummary[
  10.0,55.0,$BcMixingDefaultParameters,
  "NPoints"->20,"Progress"->False,
  AccuracyGoal->5,PrecisionGoal->5,MaxRecursion->10
];

aaMSbar = IndependentAAFinalNLOBorelMSbarSummary[
  10.0,55.0,2.0,
  $BcAlphaSMSbarDefaults,$BcMixingDefaultParameters,
  "NPoints"->20,"Progress"->False,
  AccuracyGoal->5,PrecisionGoal->5,MaxRecursion->10
];
\end{lstlisting}

\section{Mixing-angle diagnostic: what the LO and NLO numbers mean}

It is useful to keep three angle numbers separate.  They answer three
different questions.

\begin{center}
\small
\begin{tabular}{
>{\raggedright\arraybackslash}p{0.29\linewidth}
>{\raggedright\arraybackslash}p{0.35\linewidth}
>{\raggedright\arraybackslash}p{0.16\linewidth}
c}
\toprule
Name & Meaning & Perturbative content & \(\theta\)\\
\midrule
Earlier LO diagnostic &
The original central LO angle evaluated with the fixed numerical inputs
\(m_b=4.18\,{\rm GeV}\) and \(m_c=1.27\,{\rm GeV}\).  This is the number used
in the earlier LO-only discussion. &
LO only &
\(43.29^\circ\)\\
\addlinespace
Synchronized LO baseline &
The LO angle after moving all three channels to one common
\(\MSbar\)-input setup at \(\mu=4.18\,{\rm GeV}\).  This is the correct
baseline for judging the synchronized NLO diagnostic. &
LO only &
\(44.88^\circ\)\\
\addlinespace
Synchronized NLO diagnostic &
The same common \(\MSbar\)-input setup, now including the diagnostic
\(O(\as)\) corrections in \(AA\), \(AB\), and \(BB\). &
LO+\(O(\as)\) &
\(39.56^\circ\)\\
\bottomrule
\end{tabular}
\normalsize
\end{center}

The important lesson is that the NLO shift must be measured relative to the
synchronized LO baseline, not relative to the earlier fixed-input LO number:
\begin{equation}
 \Delta\theta_{\rm sync.\ diag.}
 =39.56^\circ-44.88^\circ
 =-5.32^\circ .
\end{equation}
Thus the present synchronized diagnostic says that the perturbative
\(O(\as)\) corrections lower the mixing angle by about \(5.3^\circ\) in the
common-scheme setup.

The word ``diagnostic'' is intentional.  The calculation is synchronized in a
single numerical scheme, but the \(AB\) finite coefficient is still the
current diagnostic construction rather than a closed-form coefficient derived
directly in the final \(\MSbar\) convention.  We checked the largest obvious
source of concern: the explicit tensor-current normalization
\begin{equation}
 J_B^\mu
 =
 \frac{i}{m_b+m_c}\,
 \bar b\,\sigma^{\mu\alpha}p_\alpha\gamma_5\,c .
\end{equation}
Splitting the \(AB\) pole-to-\(\MSbar\) Born-derivative mass conversion gives
\begin{align}
 \Pi^{AB}_{1,\rm mass\,conv.,full}&=0.99512,\\
 \Pi^{AB}_{1,\rm mass\,conv.,norm.\ fixed}&=0.88991,\\
 \Pi^{AB}_{1,\rm mass\,conv.,norm.\ only}&=0.10521.
\end{align}
Only about \(11\%\) of the large \(AB\) conversion comes from differentiating
the explicit \(1/(m_b+m_c)\) current normalization.  Therefore the large
\(AB\) term is not mainly a normalization artifact; it is mostly produced by
the mass dependence of the underlying \(AB\) Born spectral density.

The synchronized channel table at
\(M^2=8\,{\rm GeV}^2\), \(s_0=54\,{\rm GeV}^2\), and
\(\mu=4.18\,{\rm GeV}\) is
\begin{center}
\begin{tabular}{ccccc}
\toprule
Channel & \(\Pi_{\rm LO}\) & \(\Pi_1\) &
\((\as/\pi)\Pi_1\) & Relative shift\\
\midrule
\(AA\) & \(0.06947\) & \(-0.49883\) & \(-0.03369\) & \(-48.5\%\)\\
\(AB\) & \(-0.05573\) & \(+1.09800\) & \(+0.07415\) & \(-133.1\%\)\\
\(BB\) & \(0.06900\) & \(-0.38704\) & \(-0.02614\) & \(-37.9\%\)\\
\bottomrule
\end{tabular}
\end{center}
This table should be read as the end-of-session synchronized NLO diagnostic,
not as a final publication-grade NLO prediction.

Equivalently, the same information can be displayed as a direct
``without NLO'' versus ``with NLO'' comparison:
\begin{center}
\begin{tabular}{ccccc}
\toprule
Channel & LO & NLO shift & LO+NLO & \(K\)-factor\\
\midrule
\(AA\) & \(0.06947\) & \(-0.03369\) & \(0.03578\) & \(-0.485\)\\
\(AB\) & \(-0.05573\) & \(+0.07415\) & \(0.01842\) & \(-1.331\)\\
\(BB\) & \(0.06900\) & \(-0.02614\) & \(0.04286\) & \(-0.379\)\\
\bottomrule
\end{tabular}
\end{center}

Here ``NLO shift'' means \((\as/\pi)\Pi_1\), ``LO+NLO'' means
\(\Pi_{\rm LO}+(\as/\pi)\Pi_1\), and
\(K=(\as/\pi)\Pi_1/\Pi_{\rm LO}\).  The \(K\)-factor column is a
convergence diagnostic for each channel.  The
large value in \(AB\) should be interpreted with its sign: the LO \(AB\)
moment is negative, while the synchronized NLO change is positive, so the
relative ratio is negative and larger than unity in magnitude.

The reason such large channel corrections are plausible is the heavy-heavy
threshold region.  With the synchronized masses,
\begin{equation}
 s_{\rm th}=(\overline m_b+\overline m_c)^2=27.2898\,{\rm GeV}^2 ,
\end{equation}
and the Borel quadrature samples points very close to this threshold.  For
the first few synchronized nodes one finds
\begin{center}
\begin{tabular}{cccc}
\toprule
\(s\) & \(s-s_{\rm th}\) & \(v_{\rm rel}\) & \(\as(\mu)/v_{\rm rel}\)\\
\midrule
27.320 & 0.030 & 0.083 & 2.55\\
28.056 & 0.766 & 0.394 & 0.54\\
31.161 & 3.871 & 0.721 & 0.29\\
\bottomrule
\end{tabular}
\end{center}
Thus the lowest part of the Borel integral lies in a region where a
Coulomb-like threshold structure, schematically \(\as/v_{\rm rel}\), is not
a small number.  This explains why the channel-by-channel fixed-order NLO
diagnostics can be sizeable.  It does not, by itself, determine the sign of
each correction: the sign and final size also depend on the current
projection, the subtraction pieces, and the mass-scheme conversion.

\section{Current conclusion}

\begin{itemize}
 \item The original LO plus \(G^2\) and \(G^3\) mixing analysis is not
 invalidated by this audit.
 \item The provisional \(+3\%\) AA correction is withdrawn.  It resulted
 from reversing the epsilon-bar substitution and omitting the
 \(B_4/B_D\) normalization.
 \item The corrected independent AA coefficient passes the equal-mass
 virtual benchmark, ultraviolet and infrared cancellation, scale
 cancellation, quadrature convergence, and the massless inclusive limit.
 \item At the test point, the correction is large: approximately \(79\%\)
 in the fixed-input on-shell diagnostic and \(72\%\) in the common-scale
 \(\MSbar\) calculation at \(\mu=2\) GeV.
 \item The synchronized \(AA,AB,BB\) diagnostic gives
 \(\theta_{\rm sync.\ diag.}=39.56^\circ\), compared with the synchronized
 LO baseline \(44.88^\circ\), i.e.
 \(\Delta\theta_{\rm sync.\ diag.}=-5.32^\circ\).
 \item A publication-grade final NLO prediction still requires a fully
 independent \(AB\) finite coefficient in the same \(\MSbar\) convention.
 The immediate concern that the large \(AB\) mass-conversion term is merely
 an artifact of differentiating the tensor-current denominator
 \(1/(m_b+m_c)\) has been checked and is not supported by the split test.
\end{itemize}

\begin{thebibliography}{9}
\bibitem{Catani2002}
S.~Catani, S.~Dittmaier, M.~H.~Seymour and Z.~Trocsanyi,
Nucl.\ Phys.\ B \textbf{627} (2002) 189,
\href{https://arxiv.org/abs/hep-ph/0201036}{arXiv:hep-ph/0201036}.

\bibitem{Wang2013}
Z.-G.~Wang,
\emph{Next-to-leading order perturbative contributions in the QCD sum rules
for mesonic two-point correlation functions with unequal quark masses},
\href{https://arxiv.org/abs/1303.4146}{arXiv:1303.4146}.

\bibitem{RunDec}
K.~G.~Chetyrkin, J.~H.~K\"uhn and M.~Steinhauser,
Comput.\ Phys.\ Commun.\ \textbf{133} (2000) 43,
\href{https://arxiv.org/abs/hep-ph/0004189}{arXiv:hep-ph/0004189}.

\bibitem{ChetyrkinSteinhauser2001}
K.~G.~Chetyrkin and M.~Steinhauser,
\emph{Heavy-light current correlators at order \(\alpha_s^2\) in QCD and
HQET},
\href{https://arxiv.org/abs/hep-ph/0108017}{arXiv:hep-ph/0108017}.

\bibitem{LeeSangKim2010}
J.~Lee, W.~Sang and S.~Kim,
\emph{Relativistic corrections to the axial vector and vector currents in
the \(\bar b c\) meson system at order \(\alpha_s\)},
\href{https://arxiv.org/abs/1011.2274}{arXiv:1011.2274}.

\bibitem{MaierMarquard2011}
A.~Maier and P.~Marquard,
\emph{Low- and high-energy expansion of heavy-quark correlators at
next-to-next-to-leading order},
\href{https://arxiv.org/abs/1110.5581}{arXiv:1110.5581}.
\end{thebibliography}

\end{document}
